% FOSE2026 ライブ論文 (2ページ以内)
% 組版: platex -kanji=utf-8 fose2026; pbibtex fose2026; platex ×2; dvipdfmx fose2026.dvi
\documentclass[T,J]{fose}
\taikai{2026}

\begin{document}

% 論文のタイトル
\title{KubernetesのHPA環境下におけるコンテナ集約の適用}
\setetitle{Container Colocation of Microservices under Autoscaling}

% ★★ TODO: 著者名・所属を実際のものに差し替えてください ★★
\author{田尾瑞紀
\ejtitle{Container Colocation of Microservices\\ under Autoscaling}
\shozoku{Author Name}{和歌山大学 システム工学研究科}
{Graduate School, University}
}

\Jabstract{%
	複数サービスを1つの Pod に同居させる\textbf{部分集約}は通信コストを下げるが，
	Kubernetes では HPA が Deployment 単位で台数を決めるため，
	集約すると\textbf{スケールの粒度が粗くなり}予約資源が増える．
	本稿では Online Boutique を用いた実機実験により，集約がスループットと
	通信コストを改善する一方，この損失の大きさが\textbf{通信量}と
	\textbf{ワークロード成分}という独立な2軸で決まることを示す．
}

\maketitle \thispagestyle{empty}


\section{はじめに}

マイクロサービスを細分化すると，サービス間の呼び出しがネットワークスタックを経由するため，
通信処理のオーバヘッドが生じる．直接通信するサービスを1つの Pod に同居させ，
通信を \texttt{localhost} 経由にする\textbf{部分集約}はこれを削減する．
Wickramanayaka ら\cite{wickramanayaka2022}は Docker Swarm 環境で，
通信量の多いコンテナの適正配置・統合を行い通信量が最大58.5\%減少することを示した．

この効果を，実運用のデファクトスタンダードである Kubernetes に取り入れたい．
しかし中核機能である HPA は Deployment 単位で台数を決めるため，
集約すると\textbf{スケールの粒度が粗くなる}という問題が生じる．
本稿ではこの環境での効果を確認したうえで，相手選びに必要な判断軸を示す．


\section{部分集約とオートスケーラの干渉}

HPA は Deployment 単位で台数を決めるため，同居したコンテナは
\textbf{同じ台数でしか複製できず}，その台数は各コンテナが必要とする台数の最大値になる．
サービス $s$ の枠 (requests) を $r(s)$，単独で必要な台数を $n(s)$，
同居した Pod の台数を $N = \max_s n(s)$ とすると，余分に確保される資源 $\Delta$ は
\begin{equation}
	\Delta = \sum_{s} \bigl( N - n(s) \bigr) \times r(s)
	\label{eq:loss}
\end{equation}
となる．\textbf{損は「引きずられた側の枠」に比例}し，必要台数が揃えば $\Delta = 0$ である．
以降この $\Delta$ を\textbf{粒度損}と呼ぶ．


\section{実験環境} %追加でonline boutiqueの典拠

対象は Online Boutique\cite{onlineboutique} (11サービス，gRPC 通信) である．
MicroK8s (単一ノード，16コア) 上にデプロイし，k6を用いて負荷を与えた。

負荷は2種類の利用者を模した (表\ref{table:workload})．
\textbf{buy} は閲覧・カート追加・購入を，\textbf{view} は閲覧のみを1周とする．
\textbf{view は checkout・payment・shipping・email に一切到達しない}．
以降，1秒あたり $x$ 周のレートで与える負荷を buy($x$)・view($x$) と表す．

\begin{table}[tb]
	\centering
	\caption{2種類の利用者が到達するサービス}
	\label{table:workload}
	\small
	\doublerulesep=0.3pt
	\begin{tabular}{ll}
		\hline\hline\hline
		利用者 & 到達するサービス \\ \hline
		buy  & 全10サービス \\ \hline
		view & frontend, productcatalog, cart, \\
		     & currency, recommendation, ad \\ \hline
	\end{tabular}
\end{table}

評価指標は，粒度損として\textbf{予約枠} (稼働 Pod の requests 合計)，通信コストとして
\textbf{softirq の CPU 時間}を用いた。


\subsection{予備実験: スケールアウト後も効果は残るか}

先行研究\cite{wickramanayaka2022}が示した効果は台数1の比較である．
HPA 環境では負荷に応じて台数が増え，分離構成では複数の宛先に負荷が分散される．
そこで\textbf{台数を4に固定して HPA を切り，スケールアウト後の状態}で比較した．
構成は分離 (1 Pod 1 サービス)，frontend・recommendation・productcatalog を
同居させた \textbf{front3}，全サービスを同居させた \textbf{mega} の3つである．
全コンテナの枠を150mに揃えたので予約資源はどの構成も 6.3 コアで一致する
(3回反復，buy のみ)．

buy(150)--buy(300) ではいずれも目標レートを達成し，取りこぼしと失敗は0であった．
その条件で1周あたりの softirq は分離に対し front3 で 17--23\%，mega で 43--51\% 少なく，
\textbf{台数4でも同居させるほど通信コストが下がる}．
さらにレートを上げると分離は約360周/秒 で頭打ちになる一方，
同居は約440周/秒 まで伸び，\textbf{同じ予約資源で処理できる量が約2割多い}．


\section{本実験} %何をまとめたか記述していない

式(\ref{eq:loss})のとおり粒度損は枠 (requests) に比例するため，
枠の選び方が結果を左右する．そこで同居させるサービスの枠は実測から決めた．
各サービスを1台に固定し HPA を切った状態で負荷を6段階 (10--150周/秒) 与え，
CPU 使用量を負荷の1次式として最小二乗法で近似し，
\begin{equation}
	r = ( c_{0} + c_{1} k ) / 0.7, \quad k = 100
	\label{eq:sizing}
\end{equation}
とした．ここで $c_{0}$ はアイドル時の CPU 使用量，$c_{1}$ は1リクエストあたりの CPU 時間，
$k$ は基準負荷 [周/秒] である．
これにより，適正化したサービスは同じ負荷で目標利用率に達し，必要台数が揃うはずである．
\textbf{粒度損は同居させたサービスの枠と台数だけで決まる}ため，
適正化は同居対象に限り，他のサービスはマニフェストの値のままとした．
その上で，buy(100) を一定に保ちつつ view を0・100・200・300周/秒 の4段階で変え，
\textbf{7つの組み合わせ}を同じ条件で測った (3サイクル)．

\subsection{どの組み合わせで損が出たか}

結果を表\ref{table:pairs}に示す．
ここで\textbf{需要}とは利用率と台数の積，すなわち何台分の負荷がかかっているかを表す．
view を0から300周/秒 に増やしたとき，この需要が何倍になるかは
サービスによって 1.1--2.5倍 と幅がある．

両者を分けるのは，同居した2サービスの\textbf{伸びが揃っているかどうか}である．
伸びが近い3つでは粒度損が生じなかった．
一方，伸びが食い違う組み合わせでは損が生じ，view 比率によって
\textbf{引きずる側と引きずられる側が入れ替わった}．
ただし伸びが違っても必要台数が同じ整数に落ちれば損は出ない
(catalog+checkout)．伸びの差は損の必要条件であって十分条件ではない．

\begin{table}[tb]
	\centering
	\caption{組み合わせごとの需要の伸びと粒度損 (buy(100) 固定，view 0→300周/秒)}
	\label{table:pairs}
	\small
	\doublerulesep=0.3pt
	\setlength{\tabcolsep}{4pt}
	\begin{tabular}{lcl}
		\hline\hline\hline
		組み合わせ & 需要の伸び & 粒度損 \\ \hline
		frontend+catalog  & 2.4 / 2.1 倍 & なし \\
		frontend+reco     & 2.4 / 2.2 倍 & なし \\
		checkout+email    & 1.3 / 1.1 倍 & なし \\ \hline
		frontend+cart     & 2.4 / 1.5 倍 & \textbf{あり} \\
		catalog+checkout  & 2.1 / 1.3 倍 & ほぼなし \\
		frontend+checkout & 2.5 / 1.3 倍 & \textbf{あり} \\
		frontend+email    & 2.5 / 1.1 倍 & \textbf{あり} \\ \hline
	\end{tabular}
\end{table}

\subsection{損が枠では消せない理由}

最も比が動いた frontend+email を詳しく見る (表\ref{table:loss})．
frontend は emailservice を呼び出さないため，この組み合わせに通信削減の便益はない．

view(0) では email だけが2台を要し，frontend が引きずられて1台分が無駄になる．
view(300) では逆に frontend が3台を要し，今度は email が引きずられる．
\textbf{引きずる側と引きずられる側が入れ替わり，損が9倍変化した}．
損の比 $1077/119 = 9.05$ は枠の比と一致し，式(\ref{eq:loss})が両方向で成立している．

\begin{table}[tb]
	\centering
	\caption{view による粒度損の変化 (FE+email，buy(100) 固定，いずれも3/3で一致)}
	\label{table:loss}
	\small
	\doublerulesep=0.3pt
	\setlength{\tabcolsep}{4pt}
	\begin{tabular}{rccl}
		\hline\hline\hline
		view & 必要台数 & 需要比 & 粒度損 [mCPU] \\
		{[}周/秒] & FE/email & FE/email & (引きずられる側) \\ \hline
		0   & \textbf{1} / \textbf{2} & 0.85 & $\mathbf{+1077}$ (FE) \\ \hline
		100 & 2 / 2 & 1.24 & \textbf{0} \\ \hline
		300 & \textbf{3} / 2 & 1.90 & $\mathbf{+119}$ (email) \\ \hline
	\end{tabular}
\end{table}

view(100) では両者とも2台となり損は消える．
しかしこれは一点にすぎない．view を0から300に増やすと
\textbf{frontend の需要は2.5倍}になるのに対し\textbf{email は1.1倍}にとどまり，
両者の\textbf{需要比は2.2倍変化する}．
枠は各サービスに1つしか与えられないのに対し需要比は動くため，
\textbf{全帯で損をゼロにする枠は存在しない}．
対照として productcatalog と checkout を同居させた構成では，
需要比がほとんど動かず粒度損は12点中10点でゼロであった．

\section{考察}

便益は通信量で決まり，枠を3倍変えても softirq の削減率は 16.4\%・16.2\% と
ほぼ変わらない．一方，損は成分で決まる．この独立な2軸から，避けるべきは
\textbf{「通信量が多く，かつ成分が異なる」組み合わせ}となる．
しかし実測では太い通信の両端はいずれも同一成分で，成分をまたぐ通信は最大 2.6\% であり，
本アプリケーションに\textbf{この組み合わせは存在しなかった}．
便益の大きい相手を選ぶ限り，粒度損の危険は踏まないことになる．


\section{今後の課題}

成分を跨ぐアプリケーションでの実行。

また本実験は単一ノードで行っており，ノードをまたぐ通信 (overlay ネットワーク) を含まない．
先行研究\cite{wickramanayaka2022}が示した効果の大部分は overlay の有無に由来するため，
マルチノード環境では便益がさらに大きくなると予想される．

\bibliographystyle{jssst}
\bibliography{fose2026}

\end{document}
